% Reviewed per-case capture mechanism for all 29 examinations (final tiers after the pullback re-windows).
% Frames: Supplementary Fig. capture_cases (four per case), Fig. failure_frames (three chosen frames per failure).
% Mechanism classes: C = sustained centred ring-textured pass (success); W = dynamic wall motion / collapse;
% A = mucosal abnormality or lesion; I = instrument in the field; G = wall contact with glare / secretions;
% B = motion blur; T = ring-poor (low-texture) mucosa; L = laryngeal-only or non-airway footage; H = lens haze.
\begin{tabular}{llrlp{8.6cm}}
\toprule
Case & Tier & Gated & Class & Why it succeeded / failed (reviewed from the frames) \\
\midrule
2-V2  & CT-validated & 25 & C & Withdrawal pass over cartilage rings, lumen centred 98\% of frames, no blur or glare. \\
20-V1 & CT-validated & 24 & C, G & Centred descent (92\%); specular band on the inferior wall under-samples it but leaves the lumen measurable. \\
50-V2 & CT-validated & 19 & C & Centred proximal trachea; drift toward the carina in the last third fails the distal stations. \\
\midrule
26-V2 & measurable & 48 & C, W & \textbf{Rescued by re-windowing onto the slow centred pullback} (f1620--1901). Whole-traversal attempt failed on prolonged glottic views with an instrument; the pullback chain breaks once, at a transient posterior-wall collapse (f1912--1918), and both sides reconstruct (48/48 and 30/48). \\
12-V1 & measurable & 45 & C & Segment begins at the glottis; after the cords a centred, ring-textured descent. \\
33-V1 & measurable & 42 & C, W & Long ringed descent to the carina; a wall-angled excursion (f$\approx$750--1050) explains the discontinuity in its calibre profile. \\
27-V1 & measurable & 41 & C, G, L & Only $\sim$200 centred frames; the rest laryngeal or over-exposed (glare 29\%): measurable by count but a frayed cloud. \\
30-V2 & measurable & 40 & L & Most of the video is inside a grey artificial tube (ETT/adapter); the airway pass is brief. \\
16-V1 & measurable & 38 & G & Alternating wall close-ups and carina views: wide funnel with distal coverage loss. \\
2-V3  & measurable & 35 & C, G & Centred stretch (f$\approx$150--450) then wall-angled views producing the curled-sheet artefact. \\
15-V2 & measurable & 34 & C, L & Starts inside an ETT; centred tracheal descent (f$\approx$450--1000) gives a long tube. \\
21-V1 & measurable & 33 & C, A & Centred initial descent over irregular (cobblestoned) mucosa; later off-centre with secretions. Clinically normal. \\
22-V2 & measurable & 30 & C, A, L & Starts inside an ETT; inflamed subglottis with secretions; centred rings late (f$>$3200). \\
25-V1 & measurable & 29 & C & Brief traversal ($\sim$300 frames), centred in its second half. \\
9-V2  & measurable & 28 & C & Centred proximal trachea, then a long off-centre approach to the carina: tube ending in a bifurcation. \\
10-V2 & measurable & 25 & C, L & Long video with glottis and ETT interludes; a centred ringed descent (f$\approx$2900--3600) is measurable. \\
17-V1 & measurable & 25 & C & \textbf{Rescued by the full descent--carina--pullback loop} (f500--1330, 831/831 frames in one model): trachea, Y-shaped carina and both bronchial stubs. The whole-traversal attempt had failed on wall noise (9/48); the pullback alone gives 28/48. \\
7-V1  & measurable & 21 & C, G & Centred stretches interrupted by secretions and glare (33\%): tube with spray artefacts. \\
\midrule
24-V2 & partial & 14 & W, A, I, L & Prolonged laryngeal manipulation; compressed crescent lumen; froth; oedematous cobblestoned arytenoids; suction in view. \\
13-V1 & partial & 13 & W & \textbf{Not rescued}: descent--carina--pullback window splits into eight interleaved rigid models (vibrating posterior wall); the largest 97-frame fragment is a carina cup plus a tracheal stub. \\
14-V1 & partial & 13 & W, G & Posterior wall bulges into a crescent lumen along the ringed segment; glare 40\%. \\
31-V1 & partial & 13 & W & Clean subglottic rings (f$\approx$433--541), then a smooth wall bulge dominates the field for most of the pass. \\
1-V3  & partial & 10 & H, I & Lens haze (low contrast) with an instrument at the cords; short traversal. \\
\midrule
10-V1 & not measurable & 8 & I, G, W & Suction catheter in the lumen for most of the pass; secretions; compliant bulging wall. \\
16-V2 & not measurable & 8 & W & \textbf{Not rescued}: scope nearly static at the carina while the lumen changes shape; SfM splits into six interleaved models (99-frame fragment, coverage 0.71). \\
23-V2 & not measurable & 4 & L & Not an airway traversal: inside the endotracheal tube (printed markings visible), drapes, operating room. \\
32-V2 & not measurable & 2 & A, W, G & Inflamed irregular mucosa, an intruding wall fold, and over-exposed wall contact: mis-registered branching geometry. \\
20-V2 & not measurable & 0 & W, T, L & Collapsing lumen (round $\to$ slit $\to$ puckered star) over pale ring-poor mucosa; mostly supraglottic views. \\
24-V1 & not measurable & 0 & A, G & No lumen: wet tissue against the lens with glare and froth; rounded mass and irregular tissue (lesion). \\
\bottomrule
\end{tabular}
